\section{Knowledge: What Outsiders Do Not Know to Ask}

\subsection{Domain Mental Models}

Skill is the ability to do something. Knowledge is a model of how a domain works: its people, incentives, constraints, failure modes, vocabulary, history, workflow, and hidden trade-offs. Both matter, but knowledge often determines which opportunities become visible.

Scott Shane's study of entrepreneurs working with the same technological invention found that their previously acquired information influenced which commercial applications they recognized.\footnote{Scott Shane, ``Prior Knowledge and the Discovery of Entrepreneurial Opportunities,'' \textit{Organization Science} 11, no. 4 (2000): 448--469.} This is an illuminating mechanism, not proof that every experienced insider will identify good opportunities. A~broader meta-analysis of prior entrepreneurial experience and firm performance found a statistically significant but small overall correlation, with important variation by type of experience and context.\footnote{``Does prior experience matter? A meta-analysis of the relationship between prior experience of entrepreneurs and firm performance,'' \textit{Journal of Small Business Management} 61, no. 4 (2021): 1851--1886.}

\textbf{Experience does not eliminate uncertainty; it changes the questions you know how to ask.}

\subsection{5 Layers of Useful Domain Knowledge}

\begin{itemize}[leftmargin=11pt,itemsep=2pt,topsep=1.5pt,parsep=0pt]
  \item \textbf{Vocabulary:} You follow the industry's technical conversations. Necessary but rarely differentiating.
  \item \textbf{Workflow:} You know what happens between a customer's request and the final deliverable, including handoffs.
  \item \textbf{Incentives:} You know who benefits, who pays, who approves, and who carries the consequences of failure.
  \item \textbf{Failure modes:} You have seen where solutions break: procurement, human behavior, regulation, or integration.
  \item \textbf{Structural shifts:} You see how new technology, regulation, or cost changes make a new model feasible.
\end{itemize}

The biggest gaps are often hidden in the middle 3. People outside a market may understand the technology but fail to notice that the buyer is not the user, that approvals take 9 months, or that the proposed efficiency would eliminate a department's discretionary budget.

\subsection{The Domain-Anomaly Notebook}

For 30 days, carry a list of moments when you think: \textit{Why is this done this way?} For each anomaly, record:

\begin{enumerate}[label=\textbf{\arabic*.},leftmargin=11pt,itemsep=1.5pt,topsep=1.5pt,parsep=0pt]
  \item What is the recurring operational workaround?
  \item Who experiences the operational cost?
  \item How frequent and expensive is it?
  \item What prevents the obvious fix?
  \item Why do existing vendors tolerate it?
  \item What changed recently that could make improvement possible?
  \item What would have to be true for a buyer to switch?
  \item What am I assuming that only an outsider (or an insider) could challenge?
\end{enumerate}

Good anomalies include expensive spreadsheets, duplicated approvals, fragile integrations, repeated human explanations, regulatory bottlenecks, unreliable quality checks, and labor-intensive tasks that customers dislike doing.

A problem is not a business opportunity simply because it is annoying. Some are too cheap to solve, too rare, or tied to incentives that make change unwelcome. But the notebook creates a higher-quality supply of potential opportunities than random brainstorming.

\subsection{The Outsider's Counterweight}

Insider knowledge is valuable partly because it teaches constraints. It can also imprison you inside conventions. Experts sometimes know too well why something has always been impossible. A new technology can invalidate those reasons.

Use 2 deliberate reviews:

\begin{itemize}[leftmargin=11pt,itemsep=1.5pt,topsep=1.5pt,parsep=0pt]
  \item \textbf{Insider review:} Ask an experienced operator why your idea would fail in practice.
  \item \textbf{Outsider review:} Ask someone competent from a different field what assumptions seem arbitrary or unnecessarily fixed.
\end{itemize}

\textbf{The most attractive opportunities sit between 2 perspectives: deep operational knowledge of reality combined with willingness to challenge inherited methods.}
